% Slide 5 -- Hypothesis & objective
\begin{frame}{Hypothesis and objective}
  \begin{block}{Hypothesis}
    A rotator cuff tear produces a \alert{focal, geometrically detectable} discontinuity of the humeral
    articular contour. If we extract clinically interpretable geometric features from MRI, we can
    \alert{localize} the lesion to a specific tendon insertion sector --- without training a new neural network.
  \end{block}

  \begin{block}{Objective}
    Build an \alert{end-to-end pipeline} that:
    \begin{enumerate}
      \item Segments the humerus across the entire axial volume from a \alert{single seed slice}.
      \item Extracts two interpretable features:
        \textbf{(1)} \alert{circular arcs} on individual slices to identify those that lie on the articular surface;
        \textbf{(2)} the approximating \alert{sphere} that contains the articular surface.
      \item Uses these two mathematical objects to place the tear in the muscular complex of the rotator cuff.
    \end{enumerate}
  \end{block}
\end{frame}

% Slide 5b -- Design principles
\begin{frame}{Design principles}
  \begin{block}{Interpretable}
    Every output is geometrically meaningful and \alert{visually verifiable} by a radiologist:
    masks, boundary CSVs, fitted circles, summary figures.
  \end{block}

  \begin{block}{Vendor-agnostic}
    Works on SIEMENS, GE and Philips data \alert{without retraining}: thresholds are geometric (mm, \% of peak), not intensity-based.
  \end{block}

  \begin{block}{Anatomy-aware}
    Pipeline parameters reflect humeral anatomy (head vs shaft asymmetry, articular surface localization), \alert{not arbitrary heuristics}.
  \end{block}
\end{frame}
